\chapter{一族集合的乘积}

\section{幂集公理}

\begin{axiom}{幂集公理}{}
	$\left(\forall X\right)\Coll_{Y}\left(Y\subseteq X\right)$.
\end{axiom}

\begin{definition}{幂集}{}
	设$X$是集合.把$\left\{Y\middle|Y\subseteq X\right\}$记作$\mathcal{P}\left(X\right)$,称之为$X$的幂集.
\end{definition}

\begin{proposition}{幂集的单调性}{}
	若集合$X,Y$满足$X\subseteq Y$,则$\mathcal{P}\left(X\right)\subseteq\mathcal{P}\left(Y\right)$.
\end{proposition}

\begin{definition}{映射的典范扩张}{}
	设$\Gamma$是集合$A$到集合$B$的对应.偏函数$\hat{\Gamma}:\mathcal{P}\left(A\right)\to\mathcal{P}\left(B\right),X\mapsto\Gamma\left(X\right)$称为$\Gamma$的典范扩张.
\end{definition}

\begin{proposition}{复合映射的典范扩张}{}
	设$\Gamma$是集合$A$到集合$B$的对应,$\Gamma^{\prime}$是集合$B$到集合$C$的对应,则$\widehat{\Gamma^{\prime}\circ\Gamma}=\hat{\Gamma}^{\prime}\circ\hat{\Gamma}$.
\end{proposition}

\begin{proposition}{典范扩张的性质}{}
	设$f:E\to F$是映射.
	\begin{enumerate}
		\item 若$f$是满射(相对的,单射),则$\hat{f}$是满射(相对的,单射).
		\item 若$s$是$f$的右逆,则$\hat{s}$是$\hat{f}$的右逆.
		\item 若$r$是$f$的左逆,则$\hat{r}$是$\hat{f}$的左逆.
	\end{enumerate}
\end{proposition}


































